\chapter{Experimental Results}
\label{ch:results}

All results in this chapter are derived from the 365-day synthetic simulation dataset (seed~=~42). No result is fabricated or estimated without supporting computation.

\section{Energy System Performance (Annual)}

\begin{table}[H]
\centering
\caption{Annual energy metrics (365 days, rule-based baseline).}
\label{tab:annual_energy}
\begin{tabular}{lr}
\toprule
\textbf{Metric} & \textbf{Value} \\
\midrule
PV Generation & 5,253 kWh (14.4 kWh/day) \\
Load Consumption & 10,873 kWh (29.8 kWh/day) \\
Grid Import & 6,367 kWh (17.4 kWh/day) \\
PV Self-Consumption & 85.5\% \\
Self-Sufficiency & 41.4\% \\
Peak Load & 16,935 W \\
Peak PV & 3,143 W \\
\bottomrule
\end{tabular}
\end{table}

\section{Battery Performance}

\begin{table}[H]
\centering
\caption{Annual battery performance.}
\begin{tabular}{lr}
\toprule
\textbf{Metric} & \textbf{Value} \\
\midrule
Total Energy Charged & 1,119 kWh \\
Total Energy Discharged & 1,012 kWh \\
Round-trip Efficiency & 90.4\% \\
Average SOC & 26.9\% \\
Approximate Full Cycles & 198/year \\
\bottomrule
\end{tabular}
\end{table}

The measured round-trip efficiency (90.4\%) is consistent with the design parameters ($\eta_{\text{ch}} \times \eta_{\text{dis}} = 0.95 \times 0.95 = 0.9025$). The low average SOC (26.9\%) indicates the battery is frequently near its minimum, suggesting the 5.12~kWh capacity is well-utilized but could benefit from a larger battery in a real deployment.

\begin{figure}[H]
    \centering
    \includegraphics[width=0.75\textwidth]{figures/soc_distribution.png}
    \caption{Battery SOC distribution over 365 days.}
    \label{fig:soc_dist}
\end{figure}

At 198 cycles per year, the LiFePO4 battery's 4,000-cycle life corresponds to approximately 20 years of expected service.

\begin{figure}[H]
    \centering
    \includegraphics[width=\textwidth]{figures/monthly_energy.png}
    \caption{Monthly energy breakdown: PV generation, load consumption, and grid import.}
    \label{fig:monthly_energy}
\end{figure}

\begin{figure}[H]
    \centering
    \includegraphics[width=\textwidth]{figures/hourly_profiles.png}
    \caption{Average hourly power profiles over 365 days.}
    \label{fig:hourly_profiles}
\end{figure}

\begin{figure}[H]
    \centering
    \includegraphics[width=\textwidth]{figures/daily_self_consumption.png}
    \caption{Daily self-consumption and self-sufficiency trends.}
    \label{fig:daily_sc}
\end{figure}

\begin{figure}[H]
    \centering
    \includegraphics[width=0.85\textwidth]{figures/load_duration_curve.png}
    \caption{Load duration curve (365 days).}
    \label{fig:load_duration}
\end{figure}

\section{Optimization Results (14 Summer Days)}

\begin{table}[H]
\centering
\caption{Optimization comparison: Baseline vs Intelligent EMS (14 summer days).}
\label{tab:optimization_results}
\begin{tabular}{lrrr}
\toprule
\textbf{Metric} & \textbf{Baseline} & \textbf{Optimized} & \textbf{Improvement} \\
\midrule
Grid Import (kWh) & 213.1 & 165.1 & $-22.5\%$ \\
Total Cost (MAD) & 255.74 & 244.80 & $-4.3\%$ \\
Peak Grid Demand (W) & 14,065 & 8,840 & $-37.1\%$ \\
Self-Consumption (\%) & 81.4 & 96.0 & $+14.6$ pp \\
CO$_2$ Emissions (kg) & 138.5 & 107.3 & $-22.5\%$ \\
\bottomrule
\end{tabular}
\end{table}

The MILP optimizer achieves a 22.5\% reduction in grid imports and increases PV self-consumption from 81.4\% to 96.0\%. Peak grid demand is reduced by 37.1\% through better battery scheduling. The cost reduction is more modest (4.3\%) because the flat tariff rate limits the economic benefit of load shifting; a stronger time-of-use tariff would amplify this.

\begin{figure}[H]
    \centering
    \includegraphics[width=0.85\textwidth]{figures/optimization_comparison.png}
    \caption{Optimization comparison: baseline vs intelligent EMS across key metrics.}
    \label{fig:opt_comparison}
\end{figure}

\section{Forecasting Results Summary}

The best consumption model (XGBoost) achieves:
\begin{itemize}[nosep]
    \item Test MAE: 38.0~W
    \item Test RMSE: 73.1~W
    \item Test MAPE: 3.3\%
    \item Test $R^2$: 0.9941
\end{itemize}

This represents a 92\% reduction in RMSE compared to the naive baseline (882~W $\to$ 73~W), validating the use of ML for consumption forecasting.
